% Generated draft: requires mathematical and source review.
\documentclass[11pt]{article}
\usepackage[margin=25mm]{geometry}
\usepackage{fontspec}
\setmainfont{Latin Modern Roman}
\usepackage{amsmath,amssymb,mathtools,unicode-math}
\setmathfont{Latin Modern Math}
\usepackage{enumitem,needspace,xurl,hyperref}
\setlength{\parindent}{0pt}
\setlength{\parskip}{5pt}
\setlength{\emergencystretch}{3em}
\newcommand{\N}{\mathbb{N}}
\newcommand{\Z}{\mathbb{Z}}
\newcommand{\Q}{\mathbb{Q}}
\newcommand{\R}{\mathbb{R}}
\newcommand{\C}{\mathbb{C}}
\newcommand{\F}{\mathbb{F}}
\newcommand{\A}{\mathbb{A}}
\newcommand{\PP}{\mathbb{P}}
\newcommand{\E}{\mathbb{E}}
\newcommand{\eps}{\varepsilon}
\newcommand{\dd}{\,\mathrm{d}}
\begin{document}
\section*{UnsolvedMath ID: 800630}
\textbf{Source problem code:} catalogue entry
\subsection{Definitions and mathematical statement}
Equip \(\R\) with its usual topology, and give each \(X\subseteq\R\) the subspace topology. An open cover of \(X\) is a family of relatively open subsets of \(X\) whose union is \(X\). Let \(\mathcal O(X)\) denote the open covers of \(X\) that do not contain \(X\) itself as a member. A \(\gamma\)-cover of \(X\) is an infinite open cover \(\mathcal U\) such that, for each \(x\in X\), the set \(\{U\in\mathcal U:x\notin U\}\) is finite. Write \(\Gamma(X)\) for the family of these covers.

For \(\mathcal B(X)\in\{\mathcal O(X),\Gamma(X)\}\), the selection property \(\mathsf U_{\mathrm{fin}}(\mathcal O,\mathcal B)\) means that for every sequence \((\mathcal U_n)_{n\geq1}\) of open covers of \(X\), none of which has a finite subcover, there exist finite subfamilies \(\mathcal F_n\subseteq\mathcal U_n\) such that
\[
\left\{\bigcup\mathcal F_n:n\geq1\right\}\in\mathcal B(X).
\]
Thus the hypothesis that no \(\mathcal U_n\) has a finite subcover is part of the definition. The resulting collection is treated as a family of sets; in the \(\Gamma\) case it must be infinite.

The Menger property is \(\mathsf U_{\mathrm{fin}}(\mathcal O,\mathcal O)\), and the Hurewicz property is \(\mathsf U_{\mathrm{fin}}(\mathcal O,\Gamma)\).

\textbf{Problem.} Does ZFC prove that there exists a set \(X\subseteq\R\) such that
\[
X\models\mathsf U_{\mathrm{fin}}(\mathcal O,\mathcal O)
\qquad\text{and}\qquad
X\not\models\mathsf U_{\mathrm{fin}}(\mathcal O,\Gamma)?
\]
Here the source's phrase ``in ZFC'' is read as asking for existence without additional set-theoretic assumptions. The supplied record labels this problem open; that label is supplied curation, not an independently verified assessment of its present status. The abbreviated earlier references supplied with the problem are not bibliographically resolved in the record.
\subsection{Short English statement}
Does ZFC prove the existence of a subset of the real line that is Menger but not Hurewicz? Both properties concern finite selections from sequences of open covers having no finite subcover. The supplied open status and historical attribution require source and mathematical review.
\subsection{Sources}
\begin{itemize}
\item Local UnsolvedMath registry, permanent ID \(800630\), source problem code \texttt{catalogue entry}, upstream ID \(800630\), revision \texttt{d371ee917ac9647fbb06cd6fc529deec6f3dda82}. The supplied record attributes the problem to Tsaban (2003) and labels it open.
\item Tsaban, \emph{Selection Principles in Mathematics: A Milestone of Open Problems} (2003), Problem 2. Supplied source URL: \url{https://arxiv.org/abs/math/0312182}. The record reports access on 2026-10-16 and extraction by \texttt{arxiv\_latex}; these are supplied provenance details, not a claim of independent access or verification.
\item Supplied attribution and context: [coc2, Problem 3], [BuHa, Problem 1], and [LBopenprobs, Problem 1]. Full bibliographic entries for these reference labels are not supplied.
\end{itemize}
\end{document}
